% Experiment 9: Threaded Deterministic Reducer

\section{Experiment 9: Threaded Deterministic Reducer}

\subsection{Problem Statement}
The objective is to implement \texttt{thrsum} to sum integers 1 to N using multiple threads.

The tool must:
\begin{itemize}
    \item Accept inputs \verb|--threads <t>| and \verb|--n <N>|.
    \item Partition the range \verb|[1, N]| deterministically among threads.
    \item Use a mutex to safely update the shared total.
    \item Output \verb|OK: SUM <value>|.
\end{itemize}

\subsection{Source Code}
\begin{lstlisting}[language=C, caption={thrsum.c}]
// ######## Thredaed Deterministic Reducer ##########
// thrsum

#include <stdio.h>
#include <stdlib.h>
#include <string.h>
#include <pthread.h>
#include <stdint.h>

#define E_USAGE "E_USAGE"
#define E_RANGE "E_RANGE"
#define E_THREAD "E_THREAD"
#define E_MUTEX "E_MUTEX"

long long N_val = 0;
int num_threads = 0;
long long total_sum = 0;
pthread_mutex_t lock;

typedef struct {
    int id;
    long long start;
    long long end;
} ThreadArgs;

void print_error(const char *code, const char *msg) {
    fprintf(stderr, "ERROR: %s: %s\n", code, msg);
    exit(1);
}

void* worker(void* ptr) {
    ThreadArgs* args = (ThreadArgs*)ptr;
    long long local_sum = 0;
    
    for (long long i = args->start; i <= args->end; i++) {
        local_sum += i;
    }

    if (pthread_mutex_lock(&lock) != 0) print_error(E_MUTEX, "lock failed");
    total_sum += local_sum;
    if (pthread_mutex_unlock(&lock) != 0) print_error(E_MUTEX, "unlock failed");

    return NULL;
}

int main(int argc, char *argv[]) {
    for (int i = 1; i < argc; i++) {
        if (strcmp(argv[i], "--threads") == 0) {
            if (i + 1 < argc) num_threads = atoi(argv[++i]);
        } else if (strcmp(argv[i], "--n") == 0) {
            if (i + 1 < argc) N_val = atoll(argv[++i]);
        }
    }

    if (num_threads < 1 || num_threads > 32) print_error(E_RANGE, "threads must be in 1..32");
    if (N_val < 1 || N_val > 1000000) print_error(E_RANGE, "n must be in 1..1000000");

    if (pthread_mutex_init(&lock, NULL) != 0) print_error(E_MUTEX, "init failed");

    pthread_t threads[32];
    ThreadArgs args[32];

    long long chunk = N_val / num_threads;
    long long remainder = N_val % num_threads;
    long long current_start = 1;

    for (int i = 0; i < num_threads; i++) {
        args[i].id = i;
        args[i].start = current_start;
        long long size = chunk + (i < remainder ? 1 : 0);
        args[i].end = current_start + size - 1;
        current_start += size;

        if (pthread_create(&threads[i], NULL, worker, &args[i]) != 0) {
            print_error(E_THREAD, "create failed");
        }
    }

    for (int i = 0; i < num_threads; i++) {
        if (pthread_join(threads[i], NULL) != 0) {
            print_error(E_THREAD, "join failed");
        }
    }

    pthread_mutex_destroy(&lock);
    printf("OK: SUM %lld\n", total_sum);

    return 0;
}
\end{lstlisting}

\subsection{Sample Input and Output}

\textbf{Case 1:}
\begin{verbatim}
$ ./thrsum --threads 1 --n 10
OK: SUM 55
\end{verbatim}

\textbf{Case 2:}
\begin{verbatim}
$ ./thrsum --threads 4 --n 100
OK: SUM 5050
\end{verbatim}